% ============================================================================
%  СОПРОВОДИТЕЛЬНЫЙ МАТЕРИАЛ к курсовой работе
%  «Оптимизация обновления электрофоретического дисплея»
%  ПОЛНЫЕ доказательства Теорем 4.1 (структура оптимума) и 4.2 (нижняя оценка).
%  Компилируется отдельно:  xelatex supplement_proofs.tex  (дважды)
% ============================================================================
\documentclass[12pt,a4paper]{article}
\usepackage{fontspec}
\setmainfont{Times New Roman}
\usepackage{polyglossia}\setdefaultlanguage{russian}\setotherlanguage{english}
\usepackage[a4paper,margin=25mm]{geometry}
\usepackage{amsmath,amssymb,amsthm}
\usepackage{enumitem}
\theoremstyle{plain}
\newtheorem{theorem}{Теорема}
\newtheorem{lemma}{Лемма}
\newtheorem{proposition}{Утверждение}
\newtheorem{corollary}{Следствие}
\theoremstyle{definition}\newtheorem{remark}{Замечание}
\renewcommand{\proofname}{Доказательство}
\newcommand{\Ueff}{\mathcal{U}}
\newcommand{\Wcb}{\mathcal{W}_{\mathrm{cb}}}
\newcommand{\dstate}{d_{\mathrm{state}}}
\newcommand{\Keff}{K_{\mathrm{eff}}}
\newcommand{\xbar}{\bar{x}}

\title{Сопроводительный материал:\\полные доказательства Теорем~4.1 и~4.2}
\author{к курсовой работе «Оптимизация обновления электрофоретического дисплея»}
\date{}

\begin{document}
\maketitle

В основном тексте (глава~4) приведены формулировки теорем и схемы доказательств.
Здесь даны полные доказательства со всеми промежуточными леммами. Нумерация
теорем соответствует основному тексту (4.1, 4.2); вспомогательные леммы
нумеруются локально. Обозначения~--- из главы~4: состояние пикселя
$z[k]=(\xbar[k],q[k])^{\!\top}\in\mathbb R^{2}$ ($\xbar$~--- безразмерное
положение фронта частиц, $q$~--- остаточный заряд), управление
$u[k]\in\Ueff=\{V_{\mathrm{SH1}},V_{\mathrm{SH2}},V_{\mathrm{SL}},V_{\mathrm{COM}},0\}$
($|\Ueff|=5$), $V(\cdot)$~--- отображение кода в напряжение, $T_k\ge0$~---
длительность фазы~$k$ в кадрах, $\dstate=2$.


\section{Постановка и используемый принцип максимума}

Рассматривается дискретная задача оптимального управления
\begin{equation}\label{eq:problem}
  \min_{(\mathbf u,\mathbf T)\in\Wcb(\varepsilon)}\;
  J=\sum_{k=1}^{K} L_k(z[k],u[k],T_k),\qquad
  z[k+1]=F(z[k],u[k],T_k),
\end{equation}
где $\Wcb(\varepsilon)=\{(\mathbf u,\mathbf T):|\sum_k V(u[k])T_k|\le\varepsilon\}$~---
класс $\varepsilon$-зарядово-сбалансированных waveform, а
$L_k=\alpha\,E_k+\beta\,G_k+\gamma\,\tau_k$~--- взвешенный функционал
(энергия~$E_k=\tfrac12 C V(u[k])^2 f T_k$, гостинг~$G_k$, латентность
$\tau_k=T_k/f$) с весами $\alpha,\beta,\gamma\ge0$, $\alpha>0$.

\begin{proposition}[дискретный принцип максимума, Paruchuri и Chatterjee]
\label{prop:pmp}
Пусть отображение динамики~$F$ и функционал стоимости непрерывно дифференцируемы
по $(z,u)$, множество управлений произвольно (в частности, может быть конечным),
а активные ограничения удовлетворяют условию регулярности. Тогда для оптимального
процесса $\{(z^*[k],u^*[k])\}$ существуют сопряжённая последовательность
$\{\lambda[k]\}_{k=0}^{K}$, множитель Лагранжа~$\mu$ ограничения заряда и скаляр
$\nu\in\{0,1\}$, не равные одновременно нулю, такие что:
\begin{enumerate}[nosep]
  \item (сопряжённая система) $\displaystyle \lambda[k-1]=\frac{\partial H_k}{\partial z}$,
        $H_k=\langle\lambda[k+1],F(z[k],u[k],T_k)\rangle-\nu L_k-\mu\,V(u[k])T_k$;
  \item (условие максимума) $u^*[k]=\arg\max_{u\in\Ueff} H_k(\lambda[k+1],z^*[k],u,T_k)$;
  \item (трансверсальность) $\lambda[K]=\nu\,\partial_z\Phi$ для терминального члена~$\Phi$.
\end{enumerate}
\end{proposition}

\noindent Это формулировка теоремы~3.1 из~\cite{paruchuri2019_discrete_pmp} (обобщение
классического принципа Понтрягина на дискретное время с произвольным множеством
управлений). Ниже проверяются её условия и из условий~(i)--(ii) выводится структура
оптимума.


\section{Проверка условий применимости}

\begin{lemma}[гладкость и липшицевость динамики]\label{lem:lip}
Отображение $F$ непрерывно дифференцируемо и липшицево по $(z,u)$.
\end{lemma}
\begin{proof}
В overdamped-приближении (глава~3) интегрирование ОДУ за время $T_k$ при постоянном
управлении даёт явный аффинный вид
\begin{equation}\label{eq:Fexplicit}
  \xbar[k+1]=\xbar[k]+\tfrac{\mu^*}{L}V(u[k])\,T_k,\qquad
  q[k+1]=e^{-a T_k}q[k]+b\,V(u[k])\,T_k,
\end{equation}
с постоянными $a>0$, $b$. Каждая компонента~--- аффинная функция $z$ и линейная
функция $V(u)$, поэтому $F\in C^{\infty}$ по $(z,u)$. Якобиан $\partial_z F$ постоянен
(не зависит от $z$), а значит ограничен; следовательно $F$ глобально липшицево по~$z$
с константой $\|\partial_z F\|=\max(1,e^{-aT_k})\le1$. По~$u$ зависимость линейна,
тоже липшицева. \end{proof}

\begin{lemma}[гладкость функционала стоимости]\label{lem:cost}
$L_k$ непрерывно дифференцируем по $(z,u)$.
\end{lemma}
\begin{proof}
Энергетический член $E_k=\tfrac12 C V(u)^2 f T_k$~--- квадратичен по~$V(u)$
($C^{\infty}$); латентность $\tau_k=T_k/f$ от $z,u$ не зависит; гостинг $G_k$ есть
гладкая (квадратичная) функция отклонения состояния от целевого. Сумма гладких
функций гладка. \end{proof}

\begin{lemma}[условие регулярности и нормальность]\label{lem:slater}
При $\varepsilon>0$ выполнено условие Слейтера, поэтому множитель не вырожден:
$\nu=1$.
\end{lemma}
\begin{proof}
Ограничение $\Wcb(\varepsilon)$ есть $|\sum_k V(u[k])T_k|\le\varepsilon$. Waveform,
у которого все фазы пассивны ($u[k]=V_{\mathrm{COM}}=0$), даёт $\sum_k V(u[k])T_k=0$,
то есть лежит \emph{строго внутри} ограничения ($0<\varepsilon$). Значит допустимое
множество имеет непустую внутренность~--- условие Слейтера. По теореме о нормальности
(отсутствие абнормальных экстремалей при выполнении Слейтера) множитель при
функционале не обнуляется, и его можно нормировать: $\nu=1$. Этот свидетель
Слейтера нужен лишь для условия регулярности (непустота внутренности допустимого
множества), а не как кандидат в оптимум. \end{proof}


\section{Сопряжённая система и её решение}

\begin{lemma}[показательные моды сопряжённой переменной]\label{lem:modes}
Сопряжённая переменная имеет вид $\lambda[k]=\sum_{j} c_j\, r_j^{\,k}$, где $r_j$~---
собственные значения матрицы сопряжённой системы; для динамики~\eqref{eq:Fexplicit}
($\dstate=2$) это $r_1=1$ и $r_2=e^{aT}>1$, так что компоненты $\lambda[k]$ лежат в
линейной оболочке $\{1,\;r^{\,k}\}$ (а при кратном корне~--- $\{1,k\}$).
\end{lemma}
\begin{proof}
По~\ref{prop:pmp}(i) и~\eqref{eq:Fexplicit} сопряжённая система линейна с постоянной
матрицей $A=\partial_z F=\operatorname{diag}(1,e^{-aT})$:
$\lambda[k-1]=A^{\!\top}\lambda[k]-\nu\,\partial_z L_k$. Однородная часть~---
линейное разностное уравнение с постоянными коэффициентами; его общее решение есть
линейная комбинация $r_j^{\,k}$ по собственным значениям $A^{\!\top}$, то есть
$r_1=1$ (компонента положения~--- интегратор) и $r_2=e^{aT}$ (компонента заряда при
обращении времени). Частное решение от гладкого источника $\partial_z L_k$ не выводит
$\lambda$ из оболочки $\{1,k,r^{\,k}\}$. Оболочка $\{1,k,r^{\,k}\}$ берётся как
\emph{покрывающая} для обоих случаев (различные корни $\to\{1,r^{\,k}\}$; кратный
$\to\{1,k\}$) и даёт консервативную (не заниженную) оценку числа переключений. \end{proof}


\section{Свойство (а): релейность}

\begin{lemma}[релейность]\label{lem:bangbang}
Любое допустимое управление в~\eqref{eq:problem} релейно: $u^*[k]\in\Ueff$ для всех~$k$.
\end{lemma}
\begin{proof}
Множество $\Ueff$ конечно: уровни напряжения заданы аппаратно контроллером SSD1680,
промежуточные значения недостижимы. Условие максимума~\ref{prop:pmp}(ii) выбирает
$u^*[k]$ из $\Ueff$, поэтому каждое $u^*[k]\in\Ueff$ по построению: сигнал
кусочно-постоянен со значениями из конечного набора.

Заметим, что распространённый для непрерывных задач довод <<оптимум релаксации
над $\mathrm{conv}(\Ueff)$ достигается в вершинах>> здесь неприменим:
гамильтониан содержит вогнутый по напряжению член ($-\nu\alpha E_k\propto-V^2$), и его
максимум по выпуклой оболочке $\mathrm{conv}(\Ueff)$ мог бы достигаться во внутренней
точке. Релейность обеспечивается конечностью~$\Ueff$, а не релаксационным
аргументом. \end{proof}

\begin{lemma}[роль пассивного уровня]\label{lem:vcom}
Пассивный уровень $V_{\mathrm{COM}}=0$ минимизирует энергию фазы и потому реализует
фазы выдержки; такие фазы полярностно нейтральны и не порождают переключений знака
активного управления.
\end{lemma}
\begin{proof}
Энергия фазы $E_k=\tfrac12 C V(u)^2 f T_k$ обращается в нуль при $V(u)=0$, что
минимально среди $u\in\Ueff$. Если коэффициент при $u_k$ в~$H_k$ близок к нулю, то
прирост целевого функционала от активного драйва не окупается энергией, и оптимально
удержание $u_k=V_{\mathrm{COM}}$. Фаза с $V=0$ не меняет знака активного управления и
учитывается отдельно от активных сегментов. \end{proof}


\section{Свойство (б): оценка числа переключений}

\begin{lemma}[число знакоперемен показательной суммы]\label{lem:sc}
Ненулевая функция $g(k)=\sum_{j=0}^{m} a_j r_j^{\,k}$ с попарно различными $r_j>0$
имеет не более $m$ знакоперемен на целочисленной решётке.
\end{lemma}
\begin{proof}
Индукция по~$m$. \emph{База} $m=0$: $g(k)=a_0 r_0^{\,k}$ знака не меняет (как
произведение константы на положительную функцию). \emph{Шаг.} Разделим на
$r_0^{\,k}>0$ (это не меняет знаков): $\tilde g(k)=a_0+\sum_{j\ge1}a_j(r_j/r_0)^{\,k}$.
Дискретная разность $\Delta\tilde g(k)=\tilde g(k{+}1)-\tilde g(k)
=\sum_{j\ge1}a_j\bigl((r_j/r_0)-1\bigr)(r_j/r_0)^{\,k}$ есть показательная сумма уже из
$m$ слагаемых с попарно различными основаниями $r_j/r_0\ne1$ (постоянная $a_0$
уничтожается). По индукционному предположению $\Delta\tilde g$ имеет $\le m-1$
знакоперемен. По дискретному аналогу теоремы Ролля между двумя соседними точками
знакоперемены $\tilde g$ лежит хотя бы одна точка знакоперемены $\Delta\tilde g$;
поэтому число знакоперемен $\tilde g$ не более чем на единицу превосходит число
знакоперемен $\Delta\tilde g$, то есть $\le m$. Так как знаки $g$ и $\tilde g$
совпадают, $g$ имеет $\le m$ знакоперемен. \end{proof}

\begin{theorem}[структура оптимума, Теорема~4.1]\label{thm:main}
Оптимальное по энергии управление в классе $\Wcb(\varepsilon)$ является релейным
(bang-bang) с числом активных фаз $\Keff\le\dstate+2$ (завершающая балансирующая
фаза уже учтена в~$+2$). Для основной модели работы $\dstate=2$ (положение и
накопленный заряд), что даёт $\Keff\le4$.
\end{theorem}
\begin{proof}
Релейность установлена в лемме~\ref{lem:bangbang}. Полярность активного управления в
фазе~$k$ определяется знаком \emph{переключающей функции}
\[
  s(k)=\frac{\partial H_k}{\partial V}\Big|_{\text{акт.}}
      =\langle\lambda[k+1],B\rangle-\mu\,T_k,
\]
где $B=\partial F/\partial V$~--- вектор отклика динамики на единичное напряжение
(постоянный по~\eqref{eq:Fexplicit}). По лемме~\ref{lem:modes} компоненты $\lambda$
лежат в $\{1,k,r^{\,k}\}$, поэтому и $s(k)\in\operatorname{span}\{1,k,r^{\,k}\}$.

Применим первую разность: $\Delta s(k)=s(k{+}1)-s(k)\in\operatorname{span}\{1,r^{\,k}\}$
(линейный член $k$ даёт постоянную, $r^{\,k}$ переходит в $r^{\,k}(r-1)$). Это
показательная сумма из $m=1$ различных нетривиальных оснований плюс константа, то есть
вид $a_0+a_1 r^{\,k}$; по лемме~\ref{lem:sc} (с $m=1$) она имеет $\le1$ знакоперемену.
По дискретной теореме Ролля $s(k)$ имеет тогда $\le2=\dstate$ знакоперемен. Число
сегментов постоянной полярности на единицу больше числа знакоперемен, то есть
$\le\dstate+1=3$. С учётом не более чем одной завершающей фазы зарядового баланса
(лемма~\ref{lem:vcom}) получаем
\[
  \Keff\;\le\;\underbrace{\dstate}_{\text{знакоперемены }s}
            +\underbrace{1}_{\text{сегменты}}
            +\underbrace{1}_{\text{баланс}}
        \;=\;\dstate+2=4. \qedhere
\]
\end{proof}

\begin{remark}\label{rem:tight}
Оценка точна в смысле порядка: при разделённой динамике положения и заряда активной
оказывается лишь часть мод, и тогда $\Keff\le3$. Доказательство опирается лишь на
структуру задачи (аффинность динамики, линейность сопряжённой системы, конечность
$\Ueff$) и не зависит от конкретной формы зависимости контраста от waveform; поэтому
оно остаётся в силе и для калиброванной панельной модели главы~7, где число активных
фаз измеренного оптимума также не превосходит~4.
\end{remark}


\section{Доказательство Теоремы~4.2 (нижняя оценка энергии)}

\begin{theorem}[нижняя оценка энергии, Теорема~4.2]\label{thm:lb}
Пусть оптимальный waveform из $\Wcb(\varepsilon)$ обеспечивает изменение
отражательной способности $|\Delta\rho|=G^{*}$. Тогда:
\textnormal{(а)}~достижимость требует $G^{*}\le G_{\max}(\varepsilon)
=(\rho_{\max}-\rho_{\min})\mu^{*}\varepsilon/L$;
\textnormal{(б)}~соответствующая энергия ограничена снизу
$E^{*}\ge K_{\mathrm{lower}}\,G^{*}$ с
$K_{\mathrm{lower}}=\tfrac12 C f V_{\min}\,L/\bigl(\mu^{*}(\rho_{\max}-\rho_{\min})\bigr)>0$.
\end{theorem}

\begin{lemma}[геометрия отражательной способности]\label{lem:geom}
$|\Delta\rho|=(\rho_{\max}-\rho_{\min})\,|\Delta\xbar|$, где $\xbar\in[0,1]$.
\end{lemma}
\begin{proof}
Линейная карта положения в отражательную способность (глава~3)
$\rho=\rho_{\min}+(\rho_{\max}-\rho_{\min})\xbar$ даёт
$\Delta\rho=(\rho_{\max}-\rho_{\min})\Delta\xbar$; берём модуль. \end{proof}

\begin{lemma}[транспортное тождество]\label{lem:transport}
$\displaystyle \xbar[K]-\xbar[0]=\frac{\mu^*}{L}\sum_{k=1}^{K}V(u[k])\,T_k.$
\end{lemma}
\begin{proof}
Суммирование первой компоненты~\eqref{eq:Fexplicit} по $k=1,\dots,K$ телескопически
даёт $\xbar[K]-\xbar[0]=\frac{\mu^*}{L}\sum_k V(u[k])T_k$. \end{proof}

\paragraph{Утверждение (а).} По леммам~\ref{lem:geom},~\ref{lem:transport}
$G^{*}=(\rho_{\max}-\rho_{\min})\frac{\mu^*}{L}\bigl|\sum_k V(u[k])T_k\bigr|$.
В классе $\Wcb(\varepsilon)$ выполнено $|\sum_k V(u[k])T_k|\le\varepsilon$, откуда
$G^{*}\le(\rho_{\max}-\rho_{\min})\mu^*\varepsilon/L=G_{\max}(\varepsilon)$. $\qed$

\paragraph{Утверждение (б).} Обозначим необходимое транспортное воздействие
$Q^{*}=\bigl|\sum_k V(u[k])T_k\bigr|=\frac{L}{\mu^*}\,\frac{G^{*}}{\rho_{\max}-\rho_{\min}}$
(из лемм~\ref{lem:geom},~\ref{lem:transport}).

\begin{lemma}[оценка квадрата напряжения]\label{lem:vsq}
Для активных фаз $|V(u[k])|\ge V_{\min}>0$, поэтому $V(u[k])^2\ge V_{\min}\,|V(u[k])|$.
\end{lemma}
\begin{proof}
$V^2-V_{\min}|V|=|V|(|V|-V_{\min})\ge0$ при $|V|\ge V_{\min}$. \end{proof}

\noindent По формуле энергии и лемме~\ref{lem:vsq},
\[
  E^{*}=\sum_k\tfrac12 C\,V(u[k])^2 f\,T_k
       \;\ge\;\tfrac12 C f\,V_{\min}\sum_k |V(u[k])|\,T_k
       \;\ge\;\tfrac12 C f\,V_{\min}\Bigl|\sum_k V(u[k])T_k\Bigr|
       =\tfrac12 C f\,V_{\min}\,Q^{*},
\]
где последнее неравенство~--- неравенство треугольника
$\sum_k|V(u[k])|T_k\ge|\sum_k V(u[k])T_k|$. Подставляя $Q^{*}$, получаем
\[
  E^{*}\;\ge\;\tfrac12 C f V_{\min}\,\frac{L}{\mu^*(\rho_{\max}-\rho_{\min})}\,G^{*}
        \;=\;K_{\mathrm{lower}}\,G^{*}. \qquad\qed
\]

\begin{remark}[о бистабильности и значении константы]\label{rem:bist}
Транспортное тождество (лемма~\ref{lem:transport}) выведено в линейном
overdamped-приближении. Эксперимент (глава~7) показывает бистабильность панели:
итоговое состояние пикселя задаёт последняя фаза драйва, а энергия растёт линейно по
числу кадров. Качественный вывод теоремы (энергия не убывает с ростом требуемого
контраста и ограничена снизу) сохраняется и подтверждается измеренной зависимостью
$E^{*}(C^{*})$; уточняется лишь численное значение $K_{\mathrm{lower}}$ (его не следует
получать прямой подстановкой микроскопической подвижности~$\mu^*$, см.~\S4.5).
\end{remark}

\begin{thebibliography}{9}
\bibitem{paruchuri2019_discrete_pmp}
Paruchuri~S.~T., Chatterjee~A. A discrete-time Pontryagin maximum principle
under state-action constraints. 2019.
\end{thebibliography}

\end{document}
